\chapter{Certification Exam Mapping}
\index{certification}\index{exam mapping}\index{MongoDB Certified}%
This appendix maps the MongoDB certification tracks---the \textbf{MongoDB Certified Database
Administrator (DBA) Associate} and the \textbf{MongoDB Certified Developer Associate}---onto the
chapters of this book, then lays out a study strategy and a self-assessment checklist. Exam
blueprints evolve; treat the vendor's current exam guide as authoritative for weights and use
this appendix for the \emph{where in this book} that the vendor guide does not give you
\cite{mongodbUniversityDocs}.

\section{Certification Landscape}
\begin{center}
\footnotesize
\begin{tabular}{@{}p{3.6cm}p{4.4cm}p{4.8cm}@{}}
\toprule
\textbf{Exam} & \textbf{Audience and emphasis} & \textbf{Primary book coverage} \\
\midrule
DBA Associate & Operations: deployment, replication, sharding, security, backup, monitoring & Chapters 0--1, 5--9, 12, 17--20, 23 \\
Developer Associate & Application side: CRUD, modeling, aggregation, indexes, transactions & Chapters 1--4, 9--11, 13--14, 16 \\
Architect / advanced tracks & Multi-region design, DataOps, cost, governance & Chapters 5--8, 15--24 (this book's back half) \\
\bottomrule
\end{tabular}
\end{center}

\noindent The two Associate exams overlap heavily in the middle of the book (indexes, schema,
aggregation); the DBA exam leans into Part II (replication, sharding) and the security/operations
chapters, while the Developer exam leans into Parts I and III (modeling, aggregation, search).

\section{DBA Associate: Domain-to-Chapter Map}
\begin{center}
\footnotesize
\begin{tabular}{@{}p{4.8cm}p{3.4cm}p{4.6cm}@{}}
\toprule
\textbf{Exam domain} & \textbf{Book chapters} & \textbf{Must-be-fluent items} \\
\midrule
Installation, configuration, tooling & 0--1 & Service lifecycle, \texttt{mongod.cfg}, \texttt{serverStatus} \\
CRUD mechanics and query behavior & 0, 2 & Operators, projections, cursor behavior \\
Indexing & 9, 12 & Compound order (ESR), partial/TTL/unique, \texttt{explain} \\
Replication & 5 & Elections, oplog, write/read concern, rollback \\
Sharding & 6--8 & Shard keys, balancer, zones, \texttt{mongos}/CSRS roles \\
Security & 17--19 & SCRAM vs.\ x.509, RBAC privileges, TLS, auditing \\
Backup and restore & 20, 23 & mongodump/mongorestore, PITR, oplog replay \\
Monitoring and performance & 1, 12, 23 & Profiler, \texttt{serverStatus}, queues, alerting \\
\bottomrule
\end{tabular}
\end{center}

\section{Developer Associate: Domain-to-Chapter Map}
\begin{center}
\footnotesize
\begin{tabular}{@{}p{4.8cm}p{3.4cm}p{4.6cm}@{}}
\toprule
\textbf{Exam domain} & \textbf{Book chapters} & \textbf{Must-be-fluent items} \\
\midrule
Document model and schema design & 1--3 & Embed vs.\ reference, patterns, 16 MB limit \\
CRUD and updates & 0, 2--3 & Update operators, upsert, bulk writes, find-and-modify \\
Aggregation framework & 10--11 & Stage order, \texttt{\$group}/\texttt{\$lookup}/\texttt{\$facet}, window fields \\
Indexes and performance & 9, 12 & ESR, covered queries, reading \texttt{explain} output \\
Transactions and consistency & 5, 7 & Write/read concern, single-shard vs.\ cross-shard cost \\
Drivers and application patterns & 0, 16 & Connection lifecycle, pooling, change streams, retryable writes \\
Search and vector retrieval & 13--14 & \texttt{\$search}, analyzers, \texttt{\$vectorSearch}, filters \\
\bottomrule
\end{tabular}
\end{center}

\section{Beyond the Associate Exams}
Chapters 15--16 (stream processing, CDC, serverless), 18 (Queryable Encryption), 20--22
(governance, IaC, Kubernetes operators), and 24 (synthesis) exceed the Associate blueprints and
feed the advanced architect material---and, more importantly, the senior-engineer interview loop.
Chapter 24's synthesis questions are the bridge: each one is answerable only by chaining several
chapters, which is exactly what advanced exams and architecture interviews reward.

\section{Study Strategy}
\begin{enumerate}
    \item \textbf{Sequence.} Read Parts I--III in order; the exam domains interleave and later
    chapters assume earlier vocabulary. Use Chapter 0 to keep both environments alive.
    \item \textbf{Lab first, notes second.} Every Thursday project is an exam simulator:
    reproduce it from memory before reviewing the chapter. Recognition is not recall.
    \item \textbf{Explain everything.} For any query you write during study, produce the
    \texttt{explain("executionStats")} output and narrate it: plan chosen, docs examined, why.
    \item \textbf{Drill the failure modes.} Elections (Ch.~5), hot shards (Ch.~7), rollback
    (Ch.~5), restore (Ch.~23). Exams test what breaks, not what works.
    \item \textbf{Timed evaluation.} Sit Chapter 24's timed evaluation twice: once mid-book
    (baseline), once at the end (readiness gate). Convert misses into 48-hour gap sprints.
    \item \textbf{Close with the checklists.} Run the production-readiness checklists
    (the checklists appendix) against your capstone; each unticked box is a review target.
\end{enumerate}

\section{Self-Assessment Checklist}
You are exam-ready when every statement below is true from memory, with a lab transcript as
evidence.

\begin{itemize}
    \item[$\square$] I can draw the WiredTiger write path (cache, journal, checkpoint) and predict p99 behavior under dirty-page pressure (Ch.~1).
    \item[$\square$] I can choose embed vs.\ reference from an access-pattern inventory and defend it against the 16 MB limit (Ch.~2--3).
    \item[$\square$] I can write a \texttt{\$jsonSchema} validator and state when it does \emph{not} apply (Ch.~3).
    \item[$\square$] I can explain elections, terms, majority, and why PSA + \texttt{w:majority} has a durability hole (Ch.~5).
    \item[$\square$] I can pick \texttt{w} and read concern for a stated durability target and justify rollback behavior (Ch.~5).
    \item[$\square$] I can score a shard key on the three properties, spot a jumbo chunk forming, and fix it online (Ch.~6--7).
    \item[$\square$] I can prove zone-based residency from the chunk map, not from intentions (Ch.~8).
    \item[$\square$] I can order a compound index by ESR and prove it with \texttt{explain} (Ch.~9).
    \item[$\square$] I can write \texttt{\$match}-first pipelines, bound \texttt{\$graphLookup}, and reason about \texttt{\$facet}'s 16 MB ceiling (Ch.~10).
    \item[$\square$] I can write \texttt{\$setWindowFields} with correct \texttt{sortBy} and build an idempotent \texttt{\$merge} refresh (Ch.~11).
    \item[$\square$] I can rank slow queries by total cost and read \texttt{qr/qw} queue saturation (Ch.~12).
    \item[$\square$] I can define a Search index with autocomplete and a Vector Search index with declared filter fields (Ch.~13--14).
    \item[$\square$] I can explain exactly-once \emph{effect} via checkpoints + idempotent sinks, and the resume-token/oplog-window dependency (Ch.~15--16).
    \item[$\square$] I can design least-privilege custom roles and argue SCRAM vs.\ x.509 for services (Ch.~17).
    \item[$\square$] I can state what Queryable Encryption protects, its costs, and why KMS key loss is catastrophic (Ch.~18).
    \item[$\square$] I can enumerate a zero-trust Atlas perimeter and scoped audit filters (Ch.~19).
    \item[$\square$] I can design PITR + Online Archive + federated query, and state what PITR granularity actually equals (Ch.~20).
    \item[$\square$] I can run Terraform-driven Atlas change with plan-in-CI and forward-only migrations (Ch.~21).
    \item[$\square$] I can describe what the Kubernetes operator automates versus what remains your design (Ch.~22).
    \item[$\square$] I can execute a PITR restore drill to a new cluster and write the postmortem afterward (Ch.~23).
    \item[$\square$] I can answer Chapter 24's five synthesis questions without notes, citing measured evidence from my own labs.
\end{itemize}

\begin{notebox}
The vendor publishes the exam guide; this book publishes the competence. When they disagree,
study to the vendor guide for the sitting and to this book for the career.
\end{notebox}
